%==============================================================================
\chapter{Conclusion}
\label{sec:conclusion}
%==============================================================================

This thesis investigated whether methods independent of researcher-provided corpora can detect and mitigate hallucinations in LLM-generated Causal Loop Diagrams (CLDs), motivated by the reliability requirements of using LLMs for structured scientific knowledge generation. Across three research questions, we evaluated LLM-as-a-Judge configurations, LLM-as-a-correctors, uncertainty quantification (UQ) signals, automated literature-oriented validation via Deep Research, and adaptive deployment of Deep Research based on judge uncertainty signals.

Overall, the results support a nuanced conclusion: corpus-independent detection and correction methods can work in controlled settings, but their effectiveness degrades when the target is expert-grounded causal structure. In our dataset of three health-domain CLDs, correctness-judge performance is high in a controlled synthetic hallucination setting (GT Synth), but it drops in the literature ground-truth setting (GT Lit), with only the causal mechanistic prompt beating random baselines. In expert domains, disagreement with an expert CLD should not automatically be treated as fabrication: it can also reflect scope choices, tacit knowledge, or construct mismatch between an edge claim and what is explicitly evidenced in accessible sources. As a consequence, LLM-generated CLDs are best used as candidate drafts that focus expert attention on the contested edge space, rather than as self-validating outputs.

For corpus-independent hallucination signals, generator-side UQ metrics (perplexity, minimum probability, maximum window entropy, and ``corpus-dependent'' cosine similarity) do not effectively flag hallucinations independently. They can contribute as features to supervised ensemble classifiers in-distribution, but they fail to generalise out-of-CLD, reinforcing their role as auxiliary signals for monitoring and triage rather than as standalone detectors.

For judge-guided correction, LLM-as-a-corrector consistently improves judge scores and slightly improves CLD quality in the synthetic setting, but not versus literature ground truth. This highlights a Goodhart-like risk~\citep{goodhart1984problems} in judge--corrector feedback loops: optimizing for proxy judge satisfaction without improving expert-grounded alignment.

Finally, for literature-oriented validation, we piloted a multi-agent Deep Research (DR) system as adaptive test-time compute. DR retrieves direct causal evidence for 62.4\% of false positive edges and 42.9\% of false negative edges. Single-rater human validation (\(n=50\)) suggests that at applicability threshold \(t=0.7\), 66\% of DR-evidenced edges pass the threshold and, among passing edges, 90.9\% are correct-causal; however, a preliminary second-rater overlap (\(n=13\)) indicates only fair agreement on categorical verdicts and poor agreement on applicability scores, so calibration-based extrapolations should be interpreted as exploratory and potentially rater-sensitive. Under this calibration, extrapolated ground-truth enrichment suggests substantial gains in aggregate edge F1: \(0.267 \to 0.527\) (Scenario 1: FP\(\to\)TP) and \(0.267 \to 0.582\) (Scenario 2: LLM+DR), with conservative estimates of 0.472 and 0.500. Using a mechanistic correctness judge to flag uncertain edges for adaptive DR deployment suggests an estimated 15\% accuracy-per-dollar increase.

%------------------------------------------------------------------------------
\section{Key Contributions}
\label{sec:conclusion_contributions}
%------------------------------------------------------------------------------

\begin{itemize}[noitemsep]
    \item We provide an end-to-end experimental evaluation of corpus-independent CLD generation and post-hoc validation, spanning generation, judging, correction, uncertainty signals, and evidence-oriented verification, with methods established in literature but not tested systematically in CLD context.
    \item We show that the \textit{Correctness Judge} performs strongly in controlled synthetic hallucination settings (GT Synth) but does not reliably transfer to literature ground truth (GT Lit), where only the Mechanistic prompt exceeds random baselines.
    \item We find that the \textit{Citation Judge} does not reliably discriminate valid from hallucinated edges under either GT Synth or GT Lit: while it shows substantial alignment with a human rater when provided with relevant references (with a systematic leniency bias), weak discrimination appears to stem primarily from upstream constraints---evidence type/strength and retrieval quality. Retrieval and evidence asymmetries (notably, \textit{absence of evidence} being treated as \textit{evidence of absence} for negative edges) further make citation judging better suited to permissive screening than automated filtering.
    \item We benchmark generator-side uncertainty quantification (UQ) metrics as corpus-independent hallucination signals and show that they are weak as standalone detectors and do not generalise reliably across CLDs, but can serve as auxiliary features for monitoring and triage.
    \item We identify and characterise failure modes that matter specifically for CLDs (e.g., scope-vs-factual disagreement, plausibility bias, evidence-of-absence effects, and construct/applicability mismatch).
    \item We demonstrate that judge-guided correction can optimise for judge satisfaction without improving expert-grounded alignment, highlighting a Goodhart-like risk~\citep{goodhart1984problems} in judge--corrector feedback loops.
    \item We provide evidence that literature-oriented validation can enrich the disagreement space: calibrated with human validation and under strong assumptions (relationship isolation, a single human validator, and extrapolation), Deep Research yields substantial but bounded F1 gains; future work should test enrichment performance as these assumptions are relaxed. We also show a proof-of-concept judge-triggered deployment policy that improves cost-efficiency by selectively invoking Deep Research on uncertain edges.
\end{itemize}

%------------------------------------------------------------------------------
\section{Concluding Remarks}
\label{sec:conclusion_remarks}
%------------------------------------------------------------------------------

Taken together, these findings position LLMs as scalable assistants for proposing and stress-testing causal structures, not as autonomous validators. The most promising workflow is hybrid: use LLMs to expand the edge space, use judges and UQ signals to triage uncertainty, and selectively apply evidence-oriented validation (e.g., Deep Research with an explicit applicability threshold) to focus expert attention on the contested edges that most affect the final model.

A key finding is that ``hallucination'' in this setting is not equivalent to fabrication. In our experiments, a substantial proportion of LLM-generated edges classified as false positives against expert ground truth nonetheless had retrievable direct causal evidence under Deep Research (62.4\%). This indicates that the expert--LLM disagreement space can mix multiple phenomena: genuine fabrications, expert scoping choices, and construct or operationalisation mismatch where evidence supports a related---but not identical---claim. Accordingly, disagreement with an expert CLD should not automatically be treated as error; it can also be a signal for targeted evidence review and model refinement.

This perspective positions LLMs not merely as reproducers of expert consensus, but as tools for systematically surfacing candidate relationships for expert adjudication---an increasingly valuable function as system complexity grows and the edge space ($O(|V|^2)$) exceeds the capacity for exhaustive expert deliberation. As LLM capabilities continue to advance, the methods evaluated here---particularly uncertainty-guided selective deployment combined with literature-grounded validation---provide a foundation for integrating LLM-assisted causal synthesis into rigorous research workflows. The path forward is not to replace expert judgement, but to augment it: use LLMs to handle scalable generation and evidence gathering, while reserving expert attention for construct definition, applicability assessment, and the causal modelling decisions that require domain knowledge.

The most valuable next step would be a systematic Deep Research ablation (single-agent vs.\ multi-agent, search depth, evidence filters, and trigger policies) replicated and (human) validated, across a larger, domain-diverse CLD corpus (beyond health) to quantify generalisation and isolate which components drive gains.








